\documentclass[a4paper,12pt]{article}
\usepackage[T2A]{fontenc}
\usepackage[utf8]{inputenc}
\usepackage[russian]{babel}
\usepackage{amsmath,amssymb,mathtools}
\usepackage{paratype}
\renewcommand{\familydefault}{\sfdefault}
\usepackage{icomma}
\usepackage{geometry}
\usepackage{microtype}
\usepackage{tikz}
\usetikzlibrary{calc,arrows.meta,positioning,shapes.geometric}
\usepackage{tabularx,booktabs,longtable,array}
\usepackage{enumitem}
\usepackage{parskip}
\usepackage{fancyhdr}
\usepackage{setspace}
\usepackage{xcolor}
\usepackage{hyperref}

\definecolor{accent}{HTML}{2B6F6D}
\definecolor{darkaccent}{HTML}{184B49}
\definecolor{textgray}{HTML}{30383B}
\definecolor{softgray}{HTML}{EEF3F2}

\geometry{left=19mm,right=19mm,top=19mm,bottom=20mm}
\onehalfspacing
\setlength{\parindent}{0pt}
\setlength{\parskip}{0.42em}
\setlength{\headheight}{25pt}
\setlist[itemize]{leftmargin=1.55em,itemsep=0.25em,topsep=0.25em}
\setlist[enumerate]{leftmargin=1.8em,itemsep=0.65em,topsep=0.35em}
\hypersetup{hidelinks}

\pagestyle{fancy}
\fancyhf{}
\lhead{\small\color{textgray}Финансовые задачи ЕГЭ}
\rhead{\small\color{textgray}Тимофей Васильченко}
\cfoot{\small\color{textgray}\thepage}
\renewcommand{\headrulewidth}{0.3pt}
\renewcommand{\headrule}{\hbox to\headwidth{\color{accent}\leaders\hrule height \headrulewidth\hfill}}

\newcommand{\sectiontitle}[1]{%
  \vspace{0.25em}
  {\Large\bfseries\color{darkaccent}#1}\par
  \vspace{0.12em}
  {\color{accent}\rule{\linewidth}{0.6pt}}\par
  \vspace{0.2em}
}
\newcommand{\checkitem}{\item[$\square$]}
\newcommand{\thrub}{\text{ тыс. руб.}}
\newcommand{\mlnrub}{\text{ млн руб.}}
\newcommand{\rub}{\text{ руб.}}

\begin{document}

\begin{titlepage}
\thispagestyle{empty}
\centering
\vspace*{17mm}
{\Large\color{accent}\bfseries Чек-лист\par}
\vspace{5mm}
{\Huge\bfseries\color{darkaccent}Финансовые задачи ЕГЭ\par}
\vspace{3mm}
{\LARGE\bfseries\color{textgray}Вклады, сложные проценты и кредитная таблица\par}
\vspace{15mm}
{\large\color{textgray}Профильная математика\par}
\vspace{18mm}
\begin{tabular}{>{\bfseries}r l}
Ученик: & Тимофей Васильченко\\[2mm]
Преподаватель: & Лёвин Артём Александрович\\[2mm]
Дата: & 6 октября 2026 г.
\end{tabular}
\vfill
{\large\color{darkaccent}\bfseries
Лёвин Артём Александрович эксклюзивно для Тимофея Васильченко\par}
\vspace{10mm}
\begin{tikzpicture}[x=1cm,y=1cm]
  \draw[accent,line width=1.15pt]
    (-7.2,0)
    .. controls (-5.1,1.05) and (-3.1,-0.85) .. (-0.8,0.2)
    .. controls (1.5,1.25) and (4.1,-0.75) .. (7.2,0.15);
  \fill[accent] (-7.2,0) circle (1.35pt);
  \fill[accent] (7.2,0.15) circle (1.35pt);
\end{tikzpicture}
\vspace*{5mm}
\end{titlepage}

\sectiontitle{1. Базовая модель: деньги меняются по шагам}

В финансовой задаче сначала восстановите \textbf{порядок событий}. Один и тот же процент даст разные результаты, если пополнение, снятие или платёж происходят в разные моменты.

\begin{itemize}
\checkitem Процент \(r\%\) переводим в десятичную дробь:
\[
p=\frac{r}{100}.
\]
\checkitem Увеличение суммы на \(r\%\) означает умножение на коэффициент
\[
k=1+p=1+\frac{r}{100}.
\]
\checkitem Начисленные проценты --- это только прибавка \(pS\), а сумма после начисления --- \(S(1+p)\).
\checkitem Если единицы уже заданы словами («тысяч рублей», «миллионов рублей»), вычисления можно вести непосредственно в этих единицах.
\checkitem Округляйте только тогда, когда это требует условие. До финального ответа сохраняйте точные значения.
\end{itemize}

\sectiontitle{2. Вклад с пополнениями}

Удобная строка таблицы:
\[
\text{сумма на начало} \to \text{операция до процентов} \to \text{проценты} \to \text{сумма на конец}.
\]

\textbf{Пример.} На вклад внесли \(120\mlnrub\). В конце каждого года банк начисляет \(15\%\). В начале второго и третьего годов вклад пополняют на \(10\mlnrub\).

\[
120\cdot1{,}15=138,
\]
\[
(138+10)\cdot1{,}15=170{,}2,
\]
\[
(170{,}2+10)\cdot1{,}15=207{,}23.
\]

Итог:
\[
\boxed{207{,}23\mlnrub}.
\]

\textbf{Проверка смысла:} пополнение в начале года входит в базу начисления процентов этого же года.

\sectiontitle{3. Неизвестная начальная сумма}

Если неизвестна исходная сумма, лучше записывать всю цепочку в буквенном виде.

\textbf{Пример.} Вклад \(x\thrub\) ежегодно увеличивается на \(20\%\). В начале второго и третьего годов вносят ещё по \(10\thrub\). К концу третьего года на счёте \(173{,}28\thrub\).

\[
\bigl((1{,}2x+10)\cdot1{,}2+10\bigr)\cdot1{,}2=173{,}28.
\]

Раскрывать всё сразу необязательно. Можно двигаться обратными действиями:
\[
173{,}28:1{,}2=144{,}4,
\]
\[
144{,}4-10=134{,}4,
\]
\[
134{,}4:1{,}2=112,
\]
\[
112-10=102,
\]
\[
102:1{,}2=85.
\]

Следовательно,
\[
\boxed{x=85\thrub}.
\]

\sectiontitle{4. Неизвестная ставка и формула сложных процентов}

Если ставка постоянна и дополнительных операций нет,
\[
S_n=S_0\left(1+\frac{r}{100}\right)^n.
\]

Полезная замена:
\[
k=1+\frac{r}{100}.
\]
Тогда \(S_n=S_0k^n\).

\textbf{Пример.} \(500\thrub\) за три года превратились в \(665{,}5\thrub\):
\[
500k^3=665{,}5,
\qquad
k^3=1{,}331=\left(\frac{11}{10}\right)^3.
\]
Отсюда
\[
k=1{,}1,
\qquad
\frac{r}{100}=0{,}1,
\qquad
\boxed{r=10\%}.
\]

Если при делении возникают неудобные десятичные дроби, переводите их в обыкновенные дроби и сокращайте множители до извлечения корня.

\sectiontitle{5. Сравнение двух вкладов}

Для постоянной ставки удобно использовать степень. Для меняющейся ставки перемножайте годовые коэффициенты.

Пусть первоначально внесено \(1000\thrub\).

\textbf{Вариант 1:} три года по \(10\%\):
\[
1000\cdot1{,}1^3=1331\thrub.
\]

\textbf{Вариант 2:} \(8\%\), затем \(10\%\), затем \(12\%\):
\[
1000\cdot1{,}08\cdot1{,}10\cdot1{,}12=1330{,}56\thrub.
\]

Разница:
\[
1331-1330{,}56=0{,}44\thrub=440\rub.
\]

\textbf{Критическая проверка:} ответ записывайте именно в тех единицах, которые требует вопрос.

\sectiontitle{6. Кредит: одна универсальная таблица}

Для схемы «начислили проценты, затем внесли платёж» используйте пять величин:

\begin{center}
\begin{tabularx}{\linewidth}{>{\bfseries}l X}
\toprule
Величина & Как связана с остальными \\
\midrule
Долг на начало периода & Переносится из остатка предыдущего периода \\
Начисленные проценты & \(pD\) \\
Общая задолженность & \(D(1+p)\) \\
Платёж & \(P=D(1+p)-D_{\text{после}}\) \\
Остаток долга & \(D_{\text{после}}=D(1+p)-P\) \\
\bottomrule
\end{tabularx}
\end{center}

Главная формула сегодняшнего занятия:
\[
\boxed{P=D_{\text{до}}(1+p)-D_{\text{после}}}.
\]

Она особенно удобна, когда условие сразу сообщает, какой процент первоначального кредита должен остаться после каждого платежа.

\sectiontitle{7. Кредит по заданным остаткам}

Пусть кредит равен \(S\), месячная ставка \(5\%\), а после последовательных платежей остаётся
\[
0{,}9S;\ 0{,}8S;\ 0{,}7S;\ 0{,}6S;\ 0{,}5S;\ 0.
\]

Тогда платежи:
\[
P_1=1{,}05S-0{,}9S=0{,}15S,
\]
\[
P_2=1{,}05\cdot0{,}9S-0{,}8S=0{,}145S,
\]
\[
P_3=0{,}14S,\qquad P_4=0{,}135S,\qquad P_5=0{,}13S,
\]
\[
P_6=1{,}05\cdot0{,}5S=0{,}525S.
\]

Сумма выплат:
\[
0{,}15+0{,}145+0{,}14+0{,}135+0{,}13+0{,}525=1{,}225.
\]
Значит, выплачено \(1{,}225S\), а переплата составляет
\[
0{,}225S=22{,}5\%\text{ от суммы кредита}.
\]

\sectiontitle{8. Частые ошибки и быстрые проверки}

\begin{itemize}
\checkitem \textbf{Процент записан вместо суммы после процентов.} \(20\%\) от \(x\) --- это \(0{,}2x\), а после начисления на счёте \(1{,}2x\).
\checkitem \textbf{Проценты считают от исходной суммы, хотя база уже изменилась.} Каждый период начинайте с актуального остатка.
\checkitem \textbf{Пополнение поставлено в неправильный момент.} Проверьте формулировки «в начале», «в конце», «до начисления», «после начисления».
\checkitem \textbf{В кредитной таблице платёж считают наоборот.} Платёж = общая задолженность минус то, что осталось.
\checkitem \textbf{Теряются единицы.} \(0{,}44\) тысячи рублей = \(440\) рублей.
\checkitem \textbf{Рано переходят к большим числам.} Сначала работайте с коэффициентами \(1{,}05\), \(1{,}1\), \(1{,}2\), долями \(0{,}9S\), \(0{,}8S\) и только потом подставляйте сумму кредита.
\end{itemize}

\clearpage
\sectiontitle{9. Тренировочный блок --- Путь А}

\textbf{10 задач.} В решении сначала фиксируйте порядок операций и единицы. В этом разделе решений нет.

\begin{spacing}{1.12}
\begin{enumerate}[itemsep=0.42em]
\item На вклад внесли \(200\thrub\). В конце каждого года банк начисляет \(12\%\). Пополнений и снятий нет. Какая сумма будет на вкладе через два года?

\item На вклад внесли \(120\mlnrub\). В конце каждого года банк начисляет \(15\%\). В начале второго и третьего годов вклад пополняют на \(10\mlnrub\). Найдите сумму к концу третьего года.

\item На вклад внесли \(100\thrub\). В конце каждого года сумма увеличивается на \(20\%\). В начале второго и третьего годов вклад пополняют на \(10\thrub\). Найдите сумму к концу третьего года.

\item Начальная сумма вклада равна \(x\thrub\). В конце каждого года начисляют \(20\%\). В начале второго и третьего годов вносят ещё по \(10\thrub\). К концу третьего года на счёте \(173{,}28\thrub\). Найдите \(x\).

\item Вклад \(500\thrub\) размещён на три года без пополнений и снятий. Годовая ставка постоянна. К концу третьего года на счёте \(665{,}5\thrub\). Найдите ставку.

\item В банк внесли \(1000\thrub\) на три года. По варианту A ставка каждый год равна \(10\%\). По варианту B ставки равны соответственно \(8\%\), \(10\%\), \(12\%\). Какой вариант выгоднее и на сколько рублей?

\item Кредит \(600\thrub\) выдан на шесть месяцев под \(5\%\) в месяц. После последовательных платежей остаток долга составляет \(90\%\), \(80\%\), \(70\%\), \(60\%\), \(50\%\), \(0\%\) от первоначальной суммы. Проценты начисляются перед каждым платежом. Найдите сумму всех платежей и процент переплаты относительно кредита.

\item Кредит \(240\thrub\) взят на три года под \(10\%\) годовых. После каждого платежа тело долга уменьшается на одну и ту же величину, а к концу третьего года долг равен нулю. Найдите сумму всех платежей.

\item Кредит \(300\thrub\) взят на три года под \(20\%\) годовых. В конце первого и второго годов заёмщик платит только начисленные проценты, поэтому тело кредита остаётся неизменным. В конце третьего года кредит погашается полностью. Найдите сумму всех платежей.

\item Кредит \(400\thrub\) выдан на три года под \(10\%\) годовых. После первого платежа остаток равен \(300\thrub\), после второго --- \(180\thrub\), после третьего --- \(0\). Проценты начисляются перед каждым платежом. Найдите сумму всех платежей и процент переплаты.
\end{enumerate}
\end{spacing}

\clearpage
\sectiontitle{10. Ответы}

\renewcommand{\arraystretch}{1.38}
\begin{center}
\begin{tabularx}{\linewidth}{>{\bfseries}c X}
\toprule
№ & Ответ \\
\midrule
1 & \(250{,}88\thrub\) \\
2 & \(207{,}23\mlnrub\) \\
3 & \(199{,}2\thrub\) \\
4 & \(85\thrub\) \\
5 & \(10\%\) \\
6 & Вариант A; разница \(440\rub\) \\
7 & \(735\thrub\); переплата \(22{,}5\%\) \\
8 & \(288\thrub\) \\
9 & \(480\thrub\) \\
10 & \(488\thrub\); переплата \(22\%\) \\
\bottomrule
\end{tabularx}
\end{center}

\vfill
{\small\color{textgray}
Самопроверка: в каждом периоде сначала определите актуальную базу для процентов, затем начислите проценты, после этого выполните платёж, пополнение или снятие в том порядке, который задан условием.
}

\end{document}
